\section{A uniform saddle expansion for all reflected exponent ratios}
\label{sec:all-ratios}

Let $f(x)=\log\Gamma(x)$ and $F(x)=-\psi(x)$ on $0<x<1$.
Both functions are positive. For real $n,m\ge0$ put
\[
 M_{n,m}=\int_0^1 f(x)^n f(1-x)^m\,dx.
\]
Reflection gives $M_{n,m}=M_{m,n}$, so it is enough to consider
$n\ge m$. The previously proved compact-ratio theorem does not provide a uniform
estimate on the entire range $n/m\ge1$: its saddle in the $x$ coordinate
approaches an endpoint. The following construction uses a coordinate
and a phase which remain uniformly suitable in that limit.

\subsection{The inverse coordinate and the one needed inherited lemma}

Write $x(y)$ for the decreasing inverse of $f$, and define, for $y>0$,
\begin{equation}\label{eq:ratio-definitions}
 q(y)=f(1-x(y)),\qquad
 d(y)=\frac{e^y}{F(x(y))},\qquad
 r(y)=-\frac{yq'(y)}{q(y)}.
\end{equation}
Thus $q$ is positive and decreasing, while $d$ is positive. Direct
differentiation gives
\begin{equation}\label{eq:ratio-J}
 r(y)=\frac{1}{J(x(y))},\qquad
 J(x)=\frac{F(x)f(1-x)}{F(1-x)f(x)}.
\end{equation}
We use the following established result from the pinned repository.

\begin{lemma}[Inherited global slope inequality]\label{lem:inherited-J}
The function $J$ is an increasing real-analytic bijection
$(0,1)\longrightarrow(0,\infty)$, and
$J'(x)/J(x)>176/225$. Consequently $r$ is an increasing
real-analytic bijection $(0,\infty)\longrightarrow(0,\infty)$,
with $r'(y)>0$.
\end{lemma}

The source is the theorem entitled ``Certified global log-gamma slope
ratio'' in the report \emph{Uniform Transitions and Integral Structures},
at the revision recorded in the provenance appendix. This is an inherited
theorem, not a new assertion of this article. Its proof treats the
endpoints analytically and proves the compact inequality with 256 exact
dyadic interval computations. The dependency is included with its source
attribution and freshly replayed in this package. To verify the consequence,
differentiate \eqref{eq:ratio-J}, using $x'(y)=-1/F(x(y))$.

We also need elementary quantitative information at infinity. Put
\[
 c=\frac{\zeta(2)}{2\gamma}-\gamma>0.
\]
The inverse-function theorem applied to
$x=e^{-y}\Gamma(1+x)$ and the convergent log-gamma series yield
\begin{align}
 x(y)&=e^{-y}-\gamma e^{-2y}+O(e^{-3y}),\notag\\
 q(y)&=\gamma e^{-y}\{1+c e^{-y}+O(e^{-2y})\},\label{eq:inverse-infinity}\\
 d(y)&=1-2\gamma e^{-y}+O(e^{-2y}),\notag\\
 r(y)&=y\{1+c e^{-y}+O(e^{-2y})\}.\notag
\end{align}
Every fixed number of derivatives may be taken in these expansions.
More precisely, the correction functions in $q(y)e^y/\gamma$ and
$d(y)$ are analytic functions of $e^{-y}$ near zero. In particular,
for each fixed $j\ge2$,
\begin{equation}\label{eq:derivative-bounds}
 (\log q)^{(j)}(y)=O_j(e^{-y}),\qquad
 d^{(j)}(y)=O_j(e^{-y}),\qquad r'(y)\longrightarrow1.
\end{equation}
The statement for $d^{(j)}$ also holds for $j=1$.
Near zero, $x(y)=1-y/\gamma+O(y^2)$, so
$q(y)=\log(\gamma/y)+o(1)$ and $d(y)\to1/\gamma$.
Consequently $d$ has positive finite global upper and lower bounds.

\subsection{The all-ratio theorem}

Define
\begin{equation}\label{eq:uniform-phase}
 \Phi_{n,m}(y)=n\log y+m\log q(y)-y.
\end{equation}
The last term incorporates the exponential part of the Jacobian; keeping
it in the phase is essential when $m$ is small. The exact change of
variables is
\begin{equation}\label{eq:exact-y-moment}
 M_{n,m}=\int_0^\infty e^{\Phi_{n,m}(y)}d(y)\,dy.
\end{equation}

\begin{theorem}[Uniform expansion without a ratio restriction]
\label{thm:all-ratio}
For every real $n>0$ and $m\ge0$ there is exactly one positive solution
$s=s(n,m)$ of
\begin{equation}\label{eq:uniform-saddle}
 n=s+m r(s).
\end{equation}
It is the unique global maximizer of \eqref{eq:uniform-phase}, and
\[
 H=-\Phi_{n,m}''(s)=\frac{1+m r'(s)}{s}>0.
\]
For each fixed integer $K\ge0$, as $n\to\infty$,
\begin{equation}\label{eq:uniform-moment-expansion}
 M_{n,m}=e^{\Phi_{n,m}(s)}d(s)\sqrt{\frac{2\pi}{H}}
 \left\{\sum_{j=0}^{K}\frac{L_j(n,m)}{n^j}
                  +O_K(n^{-K-1})\right\},
 \qquad 0\le m\le n.
\end{equation}
The error constant is independent of $m$, and each coefficient
$L_j(n,m)$ is bounded on $n\ge2$, $0\le m\le n$.
Here $L_0=1$. With derivatives evaluated at $s$,
\begin{equation}\label{eq:uniform-first}
 \frac{L_1}{n}=
 \frac{d''}{2dH}
 +\frac{d'\Phi^{(3)}}{2dH^2}
 +\frac{\Phi^{(4)}}{8H^2}
 +\frac{5(\Phi^{(3)})^2}{24H^3}.
\end{equation}
The leading factor can equivalently be written
\begin{equation}\label{eq:uniform-leading-x}
 \frac{s^n q(s)^m}{F(x(s))}
       \sqrt{\frac{2\pi s}{1+m r'(s)}}.
\end{equation}
Thus reflection supplies an expansion for every pair $n,m\ge0$ as
$\max(n,m)\to\infty$.
\end{theorem}

\begin{proof}
First,
\[
 \Phi_{n,m}'(y)=\frac{n-y-mr(y)}{y}.
\]
The function $h_m(y)=y+mr(y)$ increases strictly from zero to infinity.
The derivative of the phase is therefore positive before its unique
zero and negative afterwards. At that zero its second derivative is
$-h_m'(s)/s$, proving the existence, uniqueness and curvature assertions.

The issue is uniformity, rather than this stationary calculation.
Choose $s_0>0$ so small that $s_0<1/2$ and $r(s_0)<1/2$.
For $n\ge2$, $0\le m\le n$, the saddle equation precludes $s\le s_0$:
otherwise $n=s+mr(s)<1/2+n/2<n$. Thus $s>s_0$ uniformly.
On $[s_0/2,\infty)$, positivity, continuity and
\eqref{eq:inverse-infinity} give constants $c_0,C_0,c_1,C_1>0$ such that
\begin{equation}\label{eq:uniform-comparability}
 c_0\le r(y)/y\le C_0,\qquad c_1\le r'(y)\le C_1.
\end{equation}
It follows from the saddle equation that
\[
 n\asymp(m+1)s,\qquad H\asymp n/s^2,
\]
with absolute constants throughout the stated parameter range.

For $y,s\ge s_0/2$, the mean-value estimate
$h_m'\ge c(m+1)$ gives a global phase bound. Integrating
$\Phi'=(h_m(s)-h_m(y))/y$ on either side of $s$ yields
\begin{equation}\label{eq:uniform-gap}
 \Phi(y)-\Phi(s)
 \le-c_2n\{y/s-1-\log(y/s)\}
 \qquad(y\ge s_0/2)
\end{equation}
for one constant $c_2>0$. This calculation has the same sign on both
sides: if $y<s$, integrate the positive derivative from $y$ to $s$.
The portion $0<y<s_0/2$ is also harmless. The phase is increasing there,
so its integral is at most a constant times
$e^{\Phi(s_0/2)}$. In \eqref{eq:uniform-gap},
$(s_0/2)/s\le1/2$; hence this contribution is exponentially small
relative to $e^{\Phi(s)}s/\sqrt n$.

Now put $u=y/s$ and
\[
 \varphi(u)=\frac{\Phi(su)-\Phi(s)}n,
 \qquad b(u)=\frac{d(su)}{d(s)},
 \qquad A=-\varphi''(1)=\frac{s(1+mr'(s))}{n}.
\]
The number $A$ is bounded above and bounded away from zero. Every fixed
derivative of $\varphi$ and $b$ is bounded uniformly on $1/2\le u\le3/2$.
For example, for $j\ge2$,
\[
 \varphi^{(j)}(u)
 =(-1)^{j-1}(j-1)!u^{-j}
   +\frac{m}{n}s^j(\log q)^{(j)}(su),
\]
and the second term is bounded by \eqref{eq:derivative-bounds}, compactness
on bounded $s$, and $\sup_{s\ge s_0}s^j e^{-s/2}<\infty$.
The same argument bounds the derivatives of $b$.

These estimates give a parameter-independent neighborhood of $u=1$
on which $\varphi(u)\le-c_3(u-1)^2$. Outside that neighborhood,
\eqref{eq:uniform-gap} gives an exponentially small integrable tail.
For clarity, integrability on the unbounded part follows from the linear
growth of $u-1-\log u$ as $u\to\infty$; near zero it follows from its
growth like $-\log u$.

Here is a more explicit remainder argument. Put $u=1+v/\sqrt n$ and
first restrict $|v|\le n^\epsilon$, with a positive $\epsilon$ chosen
sufficiently small for the desired fixed order $K$. On this interval,
the cubic and all higher phase terms are uniformly small relative to
the quadratic term. Taylor-expand the amplitude and the exponential
of those higher terms, retaining all powers through $n^{-K-1/2}$.
Using, for example, phase and amplitude derivatives through order
$4K+8$ and then decreasing $\epsilon$ if necessary bounds the omitted
integrand by $C_K n^{-K-1}(1+|v|^{D_K})e^{-c_4v^2}$ for a fixed
integer $D_K$. The uniform derivative bounds established above justify
these Taylor remainders. The region $|v|>n^\epsilon$ within the common
neighborhood is smaller than every power of $n^{-1}$ by the quadratic
gap; the preceding global bound handles the rest. All retained odd
powers integrate to zero after extension to the whole real line,
whose error is again smaller than every power. This proves the stated
uniform $O_K(n^{-K-1})$ normalized remainder. Boundedness of the
coefficient functions follows from the derivative bounds and the
lower bound on $A$.

For an explicit finite prescription, set
$\kappa_j=s^j\Phi^{(j)}(s)/n$ and
$b_j=s^j d^{(j)}(s)/d(s)$. If $Z$ is centered Gaussian of variance
$A^{-1}$, then $L_j$ is the coefficient of $t^{2j}$ in
\begin{equation}\label{eq:uniform-coefficient-rule}
 \mathbb E\left[
 \left(\sum_{r\ge0}\frac{b_r Z^r t^r}{r!}\right)
 \exp\left\{\sum_{r\ge3}\frac{\kappa_rZ^r t^{r-2}}{r!}\right\}
 \right].
\end{equation}
Each coefficient uses finitely many terms. The degree-two coefficient
gives \eqref{eq:uniform-first}; substitution of $d=e^s/F(x(s))$
gives \eqref{eq:uniform-leading-x}.
\end{proof}

\subsection{What uniformity does and does not imply}

The theorem includes balanced moments, fixed positive exponent ratios,
the logarithmic transition, fixed reflected exponent, and the axis $m=0$.
No lower bound on $m$ is needed. If $m$ is fixed, then
\[
 s=\frac{n}{m+1}+O_m(ne^{-n/(m+1)}),
\]
and the leading factor in \eqref{eq:uniform-moment-expansion} reduces to
the leading Stirling approximation of
$\gamma^m\Gamma(n+1)/(m+1)^{n+1}$. In the ideal limiting kernel
$q(y)=\gamma e^{-y}$, $d(y)=1$, the first coefficient is exactly
$L_1=1/12$, independently of $m$. This recovers the first Stirling
correction and checks the Jacobian and normalization.

The error in the theorem is uniform relative error of algebraic order
in the larger exponent. It does not claim the exponentially small
accuracy of the fixed-$m$ residue expansion. Likewise, the inherited
sharp least-term remainder concerns a different approximation. These
are complementary theorems with distinct error scales.

There is a useful probabilistic consequence. Give $Y$ the density
proportional to $e^{\Phi_{n,m}(y)}d(y)$. Uniformly for $0\le m\le n$,
\[
 \sqrt H\,(Y-s)\ \Longrightarrow\ N(0,1)
 \qquad(n\to\infty).
\]
Precisely, for each fixed bounded continuous $h$ and $Z\sim N(0,1)$,
\[
 \sup_{0\le m\le n}
 \left|\mathbb E h\bigl(\sqrt H(Y-s)\bigr)-\mathbb E h(Z)\right|
 \longrightarrow0.
\]
This follows from the same scaled Gaussian domination and uniform
tightness. The variance scale is $H^{-1}$;
unlike the original $x$-coordinate scale, it remains valid when the
exponent ratio tends to infinity. Higher local expansions follow from
the coefficient prescription, while a quantitative global normal-distance
bound would require an explicitly stated metric and is left for future work.
